
By naming it as \emph{case attraction}, grammarians and philologists have
assumed, since~\textcite{Lancelot1658,Lancelot1709} extended the term used
by~\textcite{Sanchez1587} for case attraction of relative pronouns,
that~\ref{gloss:attr} is some type of deviation from the general rule
of agreement\slash{}concord in AG (and Classical Latin, for that matter).
Since then, however, the cause for the variation of constructions with~\ref{gloss:attr} and
without~\ref{gloss:noattr} case attraction remained largely unexplained.
Authors have only tentatively drawn contexts in which a certain option is more
likely to occur, a handful of times assuming that such correlation between
context and resolution was causal, often due to \emph{emphasis} and
\emph{focus}, without any empirical evidence for the relation between
\emph{emphasis} and attraction.\footnote{%
	\textcite[462]{Buttmann1826},~\textcite{Hocheder1833},~\textcite[II pp.\
		24ff.]{KG},~\textcite[774]{Cooper1997},
	and~\textcite[2500--1]{Cooper2002}.
}
Other works have simply asserted its existence.\footnote{
	\textcite[203]{Goodwin1900},
	\textcite[643]{Stahl1907},
	\textcite[II p.\ 368]{Schwyzer1950},
	\textcite[p.\ 278, p.\ 440]{Smyth1956},
	\textcite[105]{Rijksbaron2002},
	\textcite{Napoli2014},
	\textcite[587]{Boas2019}.
}
% I would enroll my own works (\cite{Geraldes2020,Geraldes2021,Geraldes2025}) in
% the first group, as I mainly argued for a more precise system of correlations,
% providing quantitative evidence for such.

When it came the time for general linguists to analyse AG case agreement
patterns shown in the previous examples, their approach was rather
different, as they bundled together what a philologist would call ``normal''
case agreement and case attraction.
Since~\textcite{Lakoff1970,Andrews1971,Quicoli1982}, the works dealing with AG
case agreement patterns have been trying to provide a structural explanation for
\emph{how} non-local case concord is configured in the syntax.
Most recently, \textcite{Sevdali2013,Sevdali2013b} (drawing part of their
analysis from the works of~\cite{Tantalou2003} and~\cite{Spyropoulos2005})
argued that AG infinitive clauses may valuate their Person head locally or
externally.
For them to be locally evaluated, a ``discourse prominent'' element, not present
in the infinitive clause, must be supplied in the form of an optional accusative
\emph{pro}, which would be the case of~\ref{gloss:noattr} sentences.
Otherwise, the caseless \textsc{pro} would probe for case until it got it from
the matrix constituent which controls it, thus the case agreement
in~\ref{gloss:attr} would surface.
The cause for case attraction, under this model, would be the lack of a
``discourse prominent'' element available to be represented by an accusative
\emph{pro} in the infinitive clause.
Similarly to the philologists' ``emphasis'' explanation, however, the
``discourse prominence'' does not comes with empirical evidence or clear
definition.

The issue of ``emphasis'', ``focus'' and ``discourse prominence'' is not just a
shortcoming or unsatisfactory limit of scope in the works on case attraction.
It is also a measuring issue: there are no metrics for \emph{emphaticness} or
\emph{prominence} and the approximating techniques available for studies on
contemporary languages, namely the experimental techniques ranging from
interviews and questionnaires all the way to eye-tracking and reaction timing,
are not feasible when it comes to \emph{corpus} languages such as AG\@.



% chapter Introduction (end)  

\chapter{A stochastic approach} % (fold)

Using data from literary sources of Classical Greek, including oratory speeches,
drama, historiography and philosophical dialogues from Attic and Jonic sources,
I provide a data driven quantitative analysis of the contexts in which case
attraction is a possible agreement \slash{} concord resolution. The addition of
Jonic sources is due to the fact that it has been assumed that case attraction is
more common if not the rule in the Attic dialect
(e.g.~\cite[\emph{passim.}]{Buttmann1826} and~\cite[\emph{ad loc.}]{Cooper1997}).
The data has been collected and annotated in a combination of manual and
computational methods using the Diorisis Ancient Greek
Corpus~\parencite{TheDiorisisAncientGreekCorpus}.

A quantitative analysis of case attraction requires caution in the
methodological approach, as semantic and pragmatic features are often latent,
i.e.\ not explicitly present in the morpho-phonological level, and the use of
proxies such as word classes, constituent distance and word order may hinder the
quality of the results and the causal inference built upon them.
As such, our analysis will rely on the causality analysis (e.g.~\cite{Pearl2009})
and Bayesian modeling (e.g.~\cite{McElreath2020}), which I argue could
enhance the results of data driven linguistic research by including at the
quantitative analysis the qualitative knowledge built on the Ancient Greek
and general linguistics.
The analysis must be thus twofold. Firstly, it assess how the linguistic and
extralinguistic factors could reasonably be causally linked with case
attraction, so as to inform the statistical modeling and tell what effects are
possible to estimate from the data.
Later, a general linear model is build as to adequately estimate the direct
effects of semantic and pragmatic factors on case attraction and the
interactions between linguistic and extralinguistic variables.


% chapter A stoachastic approach (end)

\chapter{A causal model for case attraction} % (fold)

Case attraction consists in a covariance of a specific morphological feature
(case) between a controller (the oblique argument of a certain matrix verb)
and a target (the predicative or attributive expression modifying the logical
subject of the infinitive verb). Thus, it falls into the general concept of
\emph{agreement}, in Corbett's framework (\citeyear[pp.\ 4--7]{Corbett2006}).

By \emph{agreement}, we mean the general class of phenomena by which
morphological features covary across a pair of tokens with a particular
direction, i.e.:\ in which there are a controller and a target. Even there are
indeed, as~\textcite{Norris2014,Norris2017b,Polinsky2016} argue, different mechanisms
behind subject-verb agreement, DP-internal concord, and predicate case
concord --- being case attraction an instance of the last one ---, our use of
Corbett's framework will remain unaltered, as the differences observed between
DP-internal concord and subject-verb agreement in the authors defending the
contrasts are irrelevant for our analysis.



By assuming that case attraction ($A$) is a sort of \emph{non\hyp{}canonical
	agreement/concord}, it must be expected that factors known
cross\hyp{}linguistically to trigger these \emph{phenomena} also affect $A$.
Drawing a list of conditioning factors for \emph{non-canonical agreement} from
\textcite[176--205]{Corbett2006}, we should thus expect that the following
factors would change the likelihood of $A$.
\begin{inparaenum}
	\item thematic role (Θ) and animacy ($An$);
	\item grammatical relation ($GR$) and case ($K$); and
	\item communicative function ($CF$), and precedence ($Prec$).
\end{inparaenum}

% chapter A causal model for case attraction (end)



